\section{Rank-one torsion-free sheaves}

\begin{lemma}[Characterization of rank one]
\label{akcp-3-1-rank-one}
\dcref{3.1}
\source{MR0555258-compactifying-picard:page-0024}
If $X$ is geometrically integral over a field and $\II$ is coherent, then
$\II$ is rank-one torsion-free iff it is reduced (no embedded components,
length one at the generic point, and full support).  This property is preserved
and reflected by arbitrary field extension.
\end{lemma}
\begin{proof}
At the unique generic point, torsion-freeness and generic rank one are exactly
length one and absence of an embedded associated point.  These conditions are
unchanged by extension of the ground field, proving both assertions.
\end{proof}

\begin{lemma}[General hyperplane restriction]
\label{akcp-3-2-hyperplane}
\dcref{3.2}
\source{MR0555258-compactifying-picard:page-0024}
For a projective scheme over an algebraically closed field and a general
hyperplane section $Y$, restriction preserves exact sequences of coherent
modules and
\[
\Hom_X(\II,\FF)|_Y\;\xrightarrow{\sim}\;
\Hom_Y(\II|_Y,\FF|_Y).
\]
\end{lemma}
\begin{proof}
Choose $Y$ avoiding the associated points of the quotient; the snake lemma
gives exactness after restriction.  Resolve $\II$ by finite-rank locally free
modules.  Applying Hom to the resolution gives a diagram whose lower row is
exact and whose two right vertical maps are isomorphisms by restriction of a
locally free module.  The left vertical map is therefore an isomorphism.
\end{proof}

\begin{lemma}[Embedding criterion]
\label{akcp-3-3-embedding}
\dcref{3.3}
\source{MR0555258-compactifying-picard:page-0025}
\uses{akcp-3-1-rank-one,akcp-3-2-hyperplane}
For an integral projective scheme with a fixed projective embedding, a nonzero
coherent $\II$ is rank-one torsion-free iff it embeds in some $\OO_X(m)$.
For such $\II$, restriction to a general hyperplane remains rank-one
torsion-free.
\end{lemma}
\begin{proof}
An embedding in a line bundle is visibly torsion-free of generic rank one.
Conversely, Serre vanishing makes $\Hom(\II,\OO_X)(m)$ globally generated for
large $m$; a nonzero section gives a nonzero map $\II\to\OO_X(m)$, injective
because both modules are torsion-free of rank one.  Restriction follows from
the exactness in \cref{akcp-3-2-hyperplane}.
\end{proof}

\begin{proposition}[Hom between rank-one sheaves]
\label{akcp-3-4-hom-rank-one}
\dcref{3.4}
\source{MR0555258-compactifying-picard:page-0025}
\uses{akcp-3-2-hyperplane,akcp-3-3-embedding}
Let $X$ be integral projective of dimension $r\ge1$ over an algebraically
closed field, with very ample $\OO_X(1)$.  For rank-one torsion-free $\JJ,\FF$
write
$\chi(\JJ(n))=\sum a_i\binom{n}{i}$ and
$\chi(\FF(n))=\sum c_i\binom{n}{i}$.  Then $c_r=\deg X$.  If
\[
\chi(\FF(n))\le\chi(\JJ(n))\quad(n\gg0),
\]
every nonzero $\JJ\to\FF$ is an isomorphism; if the inequality is strict in
the opposite direction for all large $n$, $\Hom(\JJ,\FF)=0$.  Finally, for
$p>(c_{r-1}-a_{r-1})/\deg X+1$, $H=\Hom(\JJ,\FF)$ has $H(-p)$ a bounded
$(0,\ldots,0,\deg X)$-sheaf and is uniformly $m$-regular.
\end{proposition}
\begin{proof}
Choose a reduced zero-dimensional linear section in the locus where $\FF$ is
free.  Hilbert-polynomial slicing gives $c_r=\deg X$.  A nonzero map is
injective; its cokernel has nonnegative Hilbert function.  Under the first
inequality its Hilbert polynomial is nonpositive, hence zero, and Serre
generation forces the cokernel to vanish.  The strict reverse inequality gives
the second assertion.

For the regularity assertion, the leading coefficient of
$\chi(\FF(n))-\chi(\JJ(p+1+n))$ is negative by the bound on $p$, so the previous
vanishing argument gives $H^0(X,H(-p-1))=0$.  Slicing by a general integral
hyperplane and using \cref{akcp-3-2-hyperplane} proves inductively that $H(-p)$
is a bounded sheaf.  The regularity theorem for bounded sheaves supplies a
universal $m$.
\end{proof}

\begin{proposition}[Rank-one sheaves on an integral curve]
\label{akcp-3-5-curve-sheaves}
\dcref{3.5}
\source{MR0555258-compactifying-picard:page-0027}
\uses{akcp-3-4-hom-rank-one}
Let $X$ be an integral projective curve of arithmetic genus $p$ with dualizing
sheaf $\omega$.  For every integer $d$ there are line bundles with:
\begin{enumerate}
\item $h^0(L)=1$, $h^1(L)=d+2-p$ for $p-2<d<2p-2$;
\item $h^0(L)=p-1-d$, $h^1(L)=0$ for $d<p-2$;
\item $h^0(L)=0$, $h^1(L)=d+1-p$ for $p-1<d<2p-2$.
\end{enumerate}
There is also $\II=\omega\otimes L$ with $h^0(\II)=p-d$ and
$h^1(\II)=1$ for $0<d<p$.  For every rank-one torsion-free $\II$,
\[
\chi(\II)=\deg(\II)+1-p,\qquad
\chi(\II)<p-1\Rightarrow h^0(\II)=0,
\]
and $\chi(\II)>p-1$ implies either $h^1(\II)=0$ or $\II\simeq\omega$.
\end{proposition}
\begin{proof}
Start with $L=\OO_X$ at degree $2p-2$ and descend by tensoring with the
ideal of a smooth point.  The cohomology sequence of
$0\to L\otimes\mathfrak m_x\to L\to k(x)\to0$ gives the induction.  When
$h^1(L)>0$, duality supplies a nonzero map $L\to\omega$; choose $x$ where it
is surjective to force the evaluation map to have nontrivial kernel.  Starting
at degree $p-2$ and tensoring by the ideal gives the ascending construction in
(3).  For (ii), take $L$ from (a) with degree $2p-2-d$ and set
$\II=\omega\otimes L^{-1}$; duality computes both cohomology groups.
The Euler formula follows by slicing (or Riemann--Roch).  Applying the two
Hom alternatives of \cref{akcp-3-4-hom-rank-one} with $\OO_X$ and $\omega$
gives the final vanishing and the equality case.
\end{proof}
